% Options for packages loaded elsewhere
\PassOptionsToPackage{unicode}{hyperref}
\PassOptionsToPackage{hyphens}{url}
\documentclass[
  12pt,
]{article}
\usepackage{xcolor}
\usepackage[a4paper,left=3cm,right=2cm,top=3cm,bottom=2cm]{geometry}
\usepackage{amsmath,amssymb}
\setcounter{secnumdepth}{5}
\usepackage{iftex}
\ifPDFTeX
  \usepackage[T1]{fontenc}
  \usepackage[utf8]{inputenc}
  \usepackage{textcomp} % provide euro and other symbols
\else % if luatex or xetex
  \usepackage{unicode-math} % this also loads fontspec
  \defaultfontfeatures{Scale=MatchLowercase}
  \defaultfontfeatures[\rmfamily]{Ligatures=TeX,Scale=1}
\fi
\usepackage{lmodern}
\ifPDFTeX\else
  % xetex/luatex font selection
  \setmainfont[]{Times New Roman}
\fi
% Use upquote if available, for straight quotes in verbatim environments
\IfFileExists{upquote.sty}{\usepackage{upquote}}{}
\IfFileExists{microtype.sty}{% use microtype if available
  \usepackage[]{microtype}
  \UseMicrotypeSet[protrusion]{basicmath} % disable protrusion for tt fonts
}{}
\usepackage{setspace}
\makeatletter
\@ifundefined{KOMAClassName}{% if non-KOMA class
  \IfFileExists{parskip.sty}{%
    \usepackage{parskip}
  }{% else
    \setlength{\parindent}{0pt}
    \setlength{\parskip}{6pt plus 2pt minus 1pt}}
}{% if KOMA class
  \KOMAoptions{parskip=half}}
\makeatother
\usepackage{longtable,booktabs,array}
\newcounter{none} % for unnumbered tables
\usepackage{calc} % for calculating minipage widths
% Correct order of tables after \paragraph or \subparagraph
\usepackage{etoolbox}
\makeatletter
\patchcmd\longtable{\par}{\if@noskipsec\mbox{}\fi\par}{}{}
\makeatother
% Allow footnotes in longtable head/foot
\IfFileExists{footnotehyper.sty}{\usepackage{footnotehyper}}{\usepackage{footnote}}
\makesavenoteenv{longtable}
\usepackage{graphicx}
\makeatletter
\newsavebox\pandoc@box
\newcommand*\pandocbounded[1]{% scales image to fit in text height/width
  \sbox\pandoc@box{#1}%
  \Gscale@div\@tempa{\textheight}{\dimexpr\ht\pandoc@box+\dp\pandoc@box\relax}%
  \Gscale@div\@tempb{\linewidth}{\wd\pandoc@box}%
  \ifdim\@tempb\p@<\@tempa\p@\let\@tempa\@tempb\fi% select the smaller of both
  \ifdim\@tempa\p@<\p@\scalebox{\@tempa}{\usebox\pandoc@box}%
  \else\usebox{\pandoc@box}%
  \fi%
}
% Set default figure placement to htbp
\def\fps@figure{htbp}
\makeatother
\setlength{\emergencystretch}{3em} % prevent overfull lines
\providecommand{\tightlist}{%
  \setlength{\itemsep}{0pt}\setlength{\parskip}{0pt}}
\usepackage{indentfirst}
\usepackage{setspace}
\usepackage{caption}
\usepackage{float}
\usepackage{booktabs}
\usepackage{longtable}
\usepackage{array}
\usepackage{titlesec}
\usepackage{tocloft}
\setlength{\parindent}{1.25cm}
\setlength{\parskip}{0pt}
\captionsetup{font=small,labelfont=bf,justification=centering}
\titleformat{\section}{\normalfont\bfseries\uppercase}{\thesection}{1em}{}
\titleformat{\subsection}{\normalfont\bfseries}{\thesubsection}{1em}{}
\renewcommand{\contentsname}{Sumário}
\usepackage{bookmark}
\IfFileExists{xurl.sty}{\usepackage{xurl}}{} % add URL line breaks if available
\urlstyle{same}
\hypersetup{
  pdftitle={Associação entre notificações de dengue e temperatura média diária em Campos dos Goytacazes{,} RJ},
  pdfauthor={Análise local},
  hidelinks,
  pdfcreator={LaTeX via pandoc}}

\title{Associação entre notificações de dengue e temperatura média
diária em Campos dos Goytacazes, RJ}
\author{Análise local}
\date{03/07/2026}

\begin{document}
\maketitle

{
\setcounter{tocdepth}{3}
\tableofcontents
}
\setstretch{1.5}
\begin{titlepage}
\begin{center}
\vspace*{2cm}
{\bfseries\MakeUppercase{Associação entre notificações de dengue e temperatura média diária em Campos dos Goytacazes, RJ}\par}
\vspace{3cm}
Relatório técnico-analítico elaborado a partir de bases locais de notificações de dengue e de temperatura média diária.
\vfill
Campos dos Goytacazes, RJ\\
2026
\end{center}
\end{titlepage}

\newpage

\section{Resumo}\label{resumo}

Este relatório avaliou a associação entre notificações diárias de dengue
e temperatura média diária em Campos dos Goytacazes, RJ, no período
comum de 2020-01-08 a 2025-12-30. A análise utilizou exclusivamente
correlação de Spearman, escolhida por sua robustez frente à assimetria
das contagens, presença de muitos dias com zero casos e possibilidade de
associação monotônica não linear. Foram avaliadas associações diárias,
anomalias mensais, agregação mensal, agregação anual exploratória e
defasagens de temperatura de 0 a 45 dias. A base final incluiu 2.184
dias e 38.143 notificações. A correlação diária bruta foi positiva e
fraca (rho = 0,118, p \textless{} 0,001). A maior correlação em módulo
ocorreu com defasagem de 42 dias (rho = 0,189, p \textless{} 0,001).

\noindent\textbf{Palavras-chave:} dengue; temperatura; Spearman;
vigilância epidemiológica; séries temporais.

\newpage

\section{Materiais}\label{materiais}

Foram utilizados arquivos locais de notificações de dengue em formato
Excel (\texttt{DENGUE2020.xlsx} a \texttt{DENGUE2025.xlsx}) e a base
\texttt{temperatura\_diaria\_campos.csv}, derivada das estações
automáticas do INMET disponíveis no projeto local. A unidade geográfica
de análise foi o município de Campos dos Goytacazes, RJ, código IBGE
330100.

A variável de desfecho foi o número diário de notificações de dengue,
agregado pela data dt\_notific. A variável de exposição foi a
temperatura média diária em graus Celsius. O filtro municipal foi
baseado em notificacao. Não houve publicação, sincronização ou envio de
artefatos para GitHub.

\begin{longtable}[]{@{}
  >{\raggedright\arraybackslash}p{(\linewidth - 16\tabcolsep) * \real{0.0731}}
  >{\raggedright\arraybackslash}p{(\linewidth - 16\tabcolsep) * \real{0.0594}}
  >{\raggedright\arraybackslash}p{(\linewidth - 16\tabcolsep) * \real{0.1005}}
  >{\raggedright\arraybackslash}p{(\linewidth - 16\tabcolsep) * \real{0.1096}}
  >{\raggedright\arraybackslash}p{(\linewidth - 16\tabcolsep) * \real{0.1050}}
  >{\raggedright\arraybackslash}p{(\linewidth - 16\tabcolsep) * \real{0.1096}}
  >{\raggedright\arraybackslash}p{(\linewidth - 16\tabcolsep) * \real{0.1324}}
  >{\raggedright\arraybackslash}p{(\linewidth - 16\tabcolsep) * \real{0.1553}}
  >{\raggedright\arraybackslash}p{(\linewidth - 16\tabcolsep) * \real{0.1553}}@{}}
\caption{Resumo descritivo da base diária analisada.}\tabularnewline
\toprule\noalign{}
\begin{minipage}[b]{\linewidth}\raggedright
Dias analisados
\end{minipage} & \begin{minipage}[b]{\linewidth}\raggedright
Casos totais
\end{minipage} & \begin{minipage}[b]{\linewidth}\raggedright
Média diária de casos
\end{minipage} & \begin{minipage}[b]{\linewidth}\raggedright
Mediana diária de casos
\end{minipage} & \begin{minipage}[b]{\linewidth}\raggedright
Máximo diário de casos
\end{minipage} & \begin{minipage}[b]{\linewidth}\raggedright
Dias com zero casos (\%)
\end{minipage} & \begin{minipage}[b]{\linewidth}\raggedright
Temperatura média geral (°C)
\end{minipage} & \begin{minipage}[b]{\linewidth}\raggedright
Temperatura mínima observada (°C)
\end{minipage} & \begin{minipage}[b]{\linewidth}\raggedright
Temperatura máxima observada (°C)
\end{minipage} \\
\midrule\noalign{}
\endfirsthead
\toprule\noalign{}
\begin{minipage}[b]{\linewidth}\raggedright
Dias analisados
\end{minipage} & \begin{minipage}[b]{\linewidth}\raggedright
Casos totais
\end{minipage} & \begin{minipage}[b]{\linewidth}\raggedright
Média diária de casos
\end{minipage} & \begin{minipage}[b]{\linewidth}\raggedright
Mediana diária de casos
\end{minipage} & \begin{minipage}[b]{\linewidth}\raggedright
Máximo diário de casos
\end{minipage} & \begin{minipage}[b]{\linewidth}\raggedright
Dias com zero casos (\%)
\end{minipage} & \begin{minipage}[b]{\linewidth}\raggedright
Temperatura média geral (°C)
\end{minipage} & \begin{minipage}[b]{\linewidth}\raggedright
Temperatura mínima observada (°C)
\end{minipage} & \begin{minipage}[b]{\linewidth}\raggedright
Temperatura máxima observada (°C)
\end{minipage} \\
\midrule\noalign{}
\endhead
\bottomrule\noalign{}
\endlastfoot
2.184 & 38.143 & 17,46 & 1,00 & 458 & 49,2 & 23,84 & 15,29 & 30,49 \\
\end{longtable}

\section{Metodologia}\label{metodologia}

O processamento foi conduzido em R 4.6.0. As planilhas anuais de dengue
foram padronizadas com nomes de colunas em minúsculas, conversão de
datas, normalização dos códigos municipais e agregação diária. A base
meteorológica foi lida em formato CSV, com conversão da coluna de data e
das variáveis de temperatura para formato numérico. A série final
manteve apenas o intervalo comum entre as notificações e a temperatura.

A associação principal foi estimada por correlação de Spearman. Esse
coeficiente foi utilizado em todas as análises por não exigir
normalidade das variáveis, por reduzir a influência de extremos e por
capturar relações monotônicas. Para investigar atraso temporal plausível
entre exposição térmica e notificação, foram calculadas correlações
entre casos no dia e temperatura média observada no próprio dia e nos 45
dias anteriores.

As figuras foram produzidas com uma paleta própria inspirada em estilo
editorial Cell Press, com tons de azul profundo, verde-azulado, laranja
e dourado, fundo branco, grade leve, títulos curtos e p-valores
apresentados diretamente em cada gráfico. Este relatório segue os
elementos centrais das normas ABNT aplicáveis a trabalhos acadêmicos e
relatórios técnicos: folha de rosto, resumo, sumário, numeração
progressiva, margens em papel A4, fonte tamanho 12, espaçamento 1,5,
recuo de parágrafo e referências normativas.

\section{Resultados}\label{resultados}

\begin{longtable}[]{@{}
  >{\raggedright\arraybackslash}p{(\linewidth - 10\tabcolsep) * \real{0.4598}}
  >{\raggedright\arraybackslash}p{(\linewidth - 10\tabcolsep) * \real{0.1034}}
  >{\raggedleft\arraybackslash}p{(\linewidth - 10\tabcolsep) * \real{0.0575}}
  >{\raggedright\arraybackslash}p{(\linewidth - 10\tabcolsep) * \real{0.1839}}
  >{\raggedright\arraybackslash}p{(\linewidth - 10\tabcolsep) * \real{0.0920}}
  >{\raggedleft\arraybackslash}p{(\linewidth - 10\tabcolsep) * \real{0.1034}}@{}}
\caption{Correlações de Spearman estimadas no pipeline.}\tabularnewline
\toprule\noalign{}
\begin{minipage}[b]{\linewidth}\raggedright
Análise
\end{minipage} & \begin{minipage}[b]{\linewidth}\raggedright
Método
\end{minipage} & \begin{minipage}[b]{\linewidth}\raggedleft
n
\end{minipage} & \begin{minipage}[b]{\linewidth}\raggedright
Rho de Spearman
\end{minipage} & \begin{minipage}[b]{\linewidth}\raggedright
p-valor
\end{minipage} & \begin{minipage}[b]{\linewidth}\raggedleft
lag\_dias
\end{minipage} \\
\midrule\noalign{}
\endfirsthead
\toprule\noalign{}
\begin{minipage}[b]{\linewidth}\raggedright
Análise
\end{minipage} & \begin{minipage}[b]{\linewidth}\raggedright
Método
\end{minipage} & \begin{minipage}[b]{\linewidth}\raggedleft
n
\end{minipage} & \begin{minipage}[b]{\linewidth}\raggedright
Rho de Spearman
\end{minipage} & \begin{minipage}[b]{\linewidth}\raggedright
p-valor
\end{minipage} & \begin{minipage}[b]{\linewidth}\raggedleft
lag\_dias
\end{minipage} \\
\midrule\noalign{}
\endhead
\bottomrule\noalign{}
\endlastfoot
Diaria: temperatura media vs casos & spearman & 2184 & 0,118 &
\textless{} 0,001 & NA \\
Anomalias mensais: temperatura vs casos & spearman & 2184 & 0,190 &
\textless{} 0,001 & NA \\
Mensal: temperatura media vs casos & spearman & 72 & 0,155 & 0,195 &
NA \\
Anual: temperatura media vs casos & spearman & 6 & 0,829 & 0,042 & NA \\
Melhor defasagem diaria & spearman & 2142 & 0,189 & \textless{} 0,001 &
42 \\
\end{longtable}

A associação diária bruta entre temperatura média e notificações foi
positiva e fraca (rho = 0,118, p \textless{} 0,001). A análise por
anomalias mensais apresentou rho = 0,190 e p \textless{} 0,001,
sugerindo que parte da associação persiste após remoção simples do
padrão médio mensal. Na agregação mensal, a associação foi rho = 0,155 e
p 0,195.

\begin{figure}

{\centering \includegraphics[width=1\linewidth]{../data_inmet_campos/resultados_dengue_temperatura/figuras/00_painel_principal_cellpress} 

}

\caption{Painel principal das associações entre notificações de dengue e temperatura. Todos os painéis apresentam p-valor da correlação de Spearman correspondente.}\label{fig:fig-painel}
\end{figure}

\begin{figure}

{\centering \includegraphics[width=1\linewidth]{../data_inmet_campos/resultados_dengue_temperatura/figuras/01_serie_diaria_casos_temperatura} 

}

\caption{Série diária de notificações de dengue e temperatura média diária.}\label{fig:fig-serie}
\end{figure}

\begin{figure}

{\centering \includegraphics[width=1\linewidth]{../data_inmet_campos/resultados_dengue_temperatura/figuras/02_dispersao_diaria_spearman} 

}

\caption{Associação diária bruta entre temperatura média e notificações.}\label{fig:fig-dispersao}
\end{figure}

\begin{figure}

{\centering \includegraphics[width=1\linewidth]{../data_inmet_campos/resultados_dengue_temperatura/figuras/03_spearman_por_lag} 

}

\caption{Correlação de Spearman por defasagem da temperatura média diária.}\label{fig:fig-lag}
\end{figure}

\begin{figure}

{\centering \includegraphics[width=1\linewidth]{../data_inmet_campos/resultados_dengue_temperatura/figuras/04_agregado_mensal_spearman} 

}

\caption{Associação mensal entre temperatura média e notificações agregadas.}\label{fig:fig-mensal}
\end{figure}

\begin{figure}

{\centering \includegraphics[width=1\linewidth]{../data_inmet_campos/resultados_dengue_temperatura/figuras/05_agregado_anual_spearman} 

}

\caption{Associação anual exploratória entre temperatura média anual e notificações anuais.}\label{fig:fig-anual}
\end{figure}

\begin{figure}

{\centering \includegraphics[width=1\linewidth]{../data_inmet_campos/resultados_dengue_temperatura/figuras/06_anomalias_mensais_spearman} 

}

\caption{Associação entre anomalias mensais de temperatura e de notificações.}\label{fig:fig-anomalias}
\end{figure}

\begin{longtable}[]{@{}
  >{\raggedleft\arraybackslash}p{(\linewidth - 8\tabcolsep) * \real{0.0575}}
  >{\raggedright\arraybackslash}p{(\linewidth - 8\tabcolsep) * \real{0.2529}}
  >{\raggedright\arraybackslash}p{(\linewidth - 8\tabcolsep) * \real{0.2759}}
  >{\raggedright\arraybackslash}p{(\linewidth - 8\tabcolsep) * \real{0.1494}}
  >{\raggedright\arraybackslash}p{(\linewidth - 8\tabcolsep) * \real{0.2644}}@{}}
\caption{Resumo anual de notificações e temperatura
média.}\tabularnewline
\toprule\noalign{}
\begin{minipage}[b]{\linewidth}\raggedleft
Ano
\end{minipage} & \begin{minipage}[b]{\linewidth}\raggedright
Média diária de casos
\end{minipage} & \begin{minipage}[b]{\linewidth}\raggedright
Mediana diária de casos
\end{minipage} & \begin{minipage}[b]{\linewidth}\raggedright
Casos no ano
\end{minipage} & \begin{minipage}[b]{\linewidth}\raggedright
Temperatura média (°C)
\end{minipage} \\
\midrule\noalign{}
\endfirsthead
\toprule\noalign{}
\begin{minipage}[b]{\linewidth}\raggedleft
Ano
\end{minipage} & \begin{minipage}[b]{\linewidth}\raggedright
Média diária de casos
\end{minipage} & \begin{minipage}[b]{\linewidth}\raggedright
Mediana diária de casos
\end{minipage} & \begin{minipage}[b]{\linewidth}\raggedright
Casos no ano
\end{minipage} & \begin{minipage}[b]{\linewidth}\raggedright
Temperatura média (°C)
\end{minipage} \\
\midrule\noalign{}
\endhead
\bottomrule\noalign{}
\endlastfoot
2020 & 0,03 & 0,00 & 11 & 23,60 \\
2021 & 0,22 & 0,00 & 82 & 23,25 \\
2022 & 3,11 & 1,00 & 1.134 & 23,57 \\
2023 & 23,47 & 16,00 & 8.565 & 24,23 \\
2024 & 73,50 & 37,50 & 26.902 & 24,66 \\
2025 & 3,98 & 1,00 & 1.449 & 23,72 \\
\end{longtable}

\section{Discussão}\label{discussuxe3o}

Os resultados apontam associação positiva, porém fraca, entre
temperatura média diária e notificações de dengue. O maior coeficiente
foi observado em defasagens próximas a seis semanas, o que pode refletir
uma combinação de dinâmica vetorial, incubação, busca por atendimento e
fluxo de notificação. Ainda assim, a correlação isolada não permite
inferência causal, não ajusta simultaneamente chuva, umidade, sorotipos,
ações de controle, feriados, autocorrelação ou mudanças no sistema de
vigilância.

A leitura mais prudente é que a temperatura apresentou sinal estatístico
compatível com sazonalidade e dinâmica temporal da dengue, mas a
magnitude da associação foi baixa. Para inferência etiológica mais
robusta, recomenda-se etapa posterior com modelo de contagem, ajuste de
sazonalidade, tendência temporal, autocorrelação e, se houver dados
suficientes, modelo de defasagem distribuída.

\section{Conclusão}\label{conclusuxe3o}

No período analisado, as notificações de dengue em Campos dos Goytacazes
apresentaram associação positiva e fraca com a temperatura média diária.
A maior associação de Spearman ocorreu na análise defasada, com lag de
42 dias. Todos os resultados devem ser interpretados como análise
exploratória de associação, não como estimativa causal.

\section{Referências normativas}\label{referuxeancias-normativas}

ASSOCIAÇÃO BRASILEIRA DE NORMAS TÉCNICAS. NBR 14724: Informação e
documentação - Trabalhos acadêmicos - Apresentação. Rio de Janeiro:
ABNT.

ASSOCIAÇÃO BRASILEIRA DE NORMAS TÉCNICAS. NBR 6023: Informação e
documentação - Referências - Elaboração. Rio de Janeiro: ABNT.

ASSOCIAÇÃO BRASILEIRA DE NORMAS TÉCNICAS. NBR 6024: Informação e
documentação - Numeração progressiva das seções de um documento -
Apresentação. Rio de Janeiro: ABNT.

ASSOCIAÇÃO BRASILEIRA DE NORMAS TÉCNICAS. NBR 6027: Informação e
documentação - Sumário - Apresentação. Rio de Janeiro: ABNT.

ASSOCIAÇÃO BRASILEIRA DE NORMAS TÉCNICAS. NBR 6028: Informação e
documentação - Resumo, resenha e recensão - Apresentação. Rio de
Janeiro: ABNT.

ASSOCIAÇÃO BRASILEIRA DE NORMAS TÉCNICAS. NBR 10520: Informação e
documentação - Citações em documentos - Apresentação. Rio de Janeiro:
ABNT.

\end{document}
